\chapter*{Introduction Générale}
\markboth{\MakeUppercase{Introduction Générale}}{}
\addcontentsline{toc}{chapter}{Introduction Générale}
\adjustmtc
\thispagestyle{MyStyle}

En une dizaine d'années, la manière de bâtir et d'exploiter les logiciels a changé de
nature. Portés par le Cloud Computing et le DevOps, les projets ont quitté les serveurs
traditionnels au profit d'architectures \textit{cloud-native}, fondées sur la
conteneurisation, l'\textit{Infrastructure as Code} et les pipelines d'intégration et de
livraison continues (CI/CD). À la clé~: des livraisons plus rapides, une forte automatisation
et des environnements aisément reproductibles.

Cette évolution a son revers, car la complexité grimpe. Mettre une application en production
oblige à composer avec un large éventail d'outils~--- Docker, Kubernetes, Terraform,
fournisseurs cloud~--- et à trancher, étape après étape, quantité de choix techniques.
Architecture, dimensionnement des ressources, agencement des artefacts, pipelines~: le moindre
faux pas de conception retentit sur la fiabilité comme sur le coût de l'infrastructure.

L'intelligence artificielle générative et les grands modèles de langage (LLM) laissent
entrevoir une aide réelle, sans pour autant remplacer l'expertise humaine~: leur rôle est
d'assister, non de décider. Les assistants actuels restent pourtant partiels. Beaucoup se
limitent à produire des fragments de code ou des configurations isolés, sans vue d'ensemble
du projet~; rares sont ceux qui analysent un dépôt de bout en bout, proposent une architecture
cloud adaptée aux contraintes, préparent des artefacts cohérents, font valider les choix ou
corrigent finement les erreurs. Faute de mieux, l'ingénieur doit recoller lui-même des
morceaux épars, sans garantie de cohérence.

C'est à ce manque que s'attaque le présent travail, réalisé lors d'un stage de fin d'études
chez Talan Tunisie Consulting. Nous y présentons \textbf{MultiMindOps}, une plateforme
d'assistance IA destinée aux ingénieurs Cloud et DevOps. Plutôt que de livrer des fichiers
détachés les uns des autres, elle accompagne l'ingénieur du début à la fin~: analyse
automatique du projet, propositions d'architectures cloud, estimation des ressources,
génération des principaux artefacts DevOps et validation assistée des décisions. Le principe
directeur en est le \textit{Human-in-the-Loop}~: l'ingénieur garde la maîtrise des
orientations stratégiques et approuve chaque étape. La sécurité, quant à elle, n'y forme pas
une brique isolée mais une exigence de qualité transversale, présente autant dans les
artefacts produits que dans l'architecture de la plateforme.

\bigskip

\noindent La suite du document s'organise ainsi :

\begin{itemize}
    \item \textbf{Chapitre 1 --- Cadre général du projet~:} il campe l'organisme d'accueil,
    précise le contexte et la problématique, passe en revue les outils existants, avant
    d'exposer la solution retenue et la méthodologie de travail.

    \item \textbf{Chapitre 2 --- Concepts fondamentaux~:} il réunit l'état de l'art des
    notions mobilisées~: intelligence artificielle générative et LLM, systèmes multi-agents,
    RAG, Infrastructure as Code, conteneurisation et orchestration cloud.

    \item \textbf{Chapitre 3 --- Analyse et conception~:} il traite l'analyse des besoins,
    la modélisation de l'architecture d'ensemble, les diagrammes UML et la conception fine
    des agents et des workflows.

    \item \textbf{Chapitre 4 --- Réalisation et mise en œuvre~:} il rend compte de
    l'implémentation, du déploiement des agents IA, de la génération des artefacts DevOps,
    des interfaces utilisateur et d'une démonstration des fonctionnalités.
\end{itemize}

\bigskip

\noindent Une conclusion générale clôt l'ensemble~: elle fait le bilan des apports de ce
travail et esquisse les pistes d'évolution envisagées pour la plateforme.
